\documentclass[10pt,a4paper]{amsart}
\usepackage[margin=19mm]{geometry}
\usepackage{amsmath,amssymb,mathtools}
\usepackage[scr=rsfs,cal=euler]{mathalfa}
\usepackage{xfrac}
\usepackage[unicode,hidelinks,bookmarks=false]{hyperref}
\newcommand{\ie}{i.e.\ }
\newcommand{\cf}{cf.\ }
\newcommand{\resp}{respectively\ }
\newcommand{\Subsection}{\subsection}
\providecommand{\nobreakdash}{\nobreak\hbox{-}}
\setlength{\emergencystretch}{2em}
\multlinegap0pt
\usepackage{mathtools}

\usepackage{wasysym}
\usepackage{mathrsfs}

\newcommand*\ssS{\mathsf{S}}
\newcommand*\GS{\mathsf{GS}}
\newcommand*\ssN{\mathsf{N}}
\newcommand*\sfs{\mathsf{s}}
\newcommand*\sfr{\mathsf{r}}
\newcommand*\sff{\mathsf{f}}
\newcommand*\sfT{\mathsf{T}}


\newcommand*\scM{\mathscr{M}}
\newcommand*\scN{\mathscr{N}}

\newcommand*\TT{\mathbb{T}}
\newcommand*\RR{\mathbb{R}}
\newcommand*\ZZ{\mathbb{Z}}
\newcommand*\NN{\mathbb{N}}
\newcommand{\bbQ}{\mathbb{Q}}
\newcommand{\bbP}{\mathbb{P}}
\newcommand{\bbU}{\mathbb{U}}
\newcommand{\bbE}{\mathbb{E}}

\newcommand{\bfun}{\mathbf{1}}
\newcommand{\clA}{\mathcal{A}}
\newcommand{\clJ}{\mathcal{J}}
\newcommand{\clR}{\mathcal{R}}
\newcommand{\clE}{\mathcal{E}}
\newcommand{\clF}{\mathcal{F}}
\newcommand{\clT}{\mathcal{T}}



\newcommand*\st{\mathrel{:}}

\DeclareMathOperator{\supp}{supp}
\DeclareMathOperator{\locc}{loc}
\newcommand*\e{\mathrm{e}}
\newcommand{\rms}{\mathrm{s}}
\newcommand*\dif{\mathrm{d}}
\newcommand*\Left{\mathrm{L}}
\newcommand*\Right{\mathrm{R}}
\newcommand*\LR{\mathrm{LR}}

\newcommand*\eps{\varepsilon}

\renewcommand*{\to}{\mathchoice{\longrightarrow}{\rightarrow}{\rightarrow}{\rightarrow}}

\let\oldmapsto\mapsto
\renewcommand*{\mapsto}{\mathchoice{\longmapsto}{\oldmapsto}{\oldmapsto}{\oldmapsto}}

\let\oldin\in
\renewcommand{\in}{\mathchoice{\oldin}{\oldin}{\,\oldin\,}{\,\oldin\,}}
\let\oldcap\cap
\renewcommand{\cap}{\mathchoice{\oldcap}{\oldcap}{\,\oldcap\,}{\,\oldcap\,}}

\let\originalleft\left 
\let\originalright\right
\renewcommand{\left}{\mathopen{}\mathclose\bgroup\originalleft}
\renewcommand{\right}{\aftergroup\egroup\originalright}

\newcommand*{\mk}{\mkern -1mu}
\newcommand*{\Mk}{\mkern -2mu}
\newcommand*{\mK}{\mkern 1mu}
\newcommand*{\MK}{\mkern 2mu}

\hypersetup{urlcolor=purple, linkcolor=blue, citecolor=red}

\newcommand*{\relabel}{\renewcommand{\labelenumi}{(\theenumi)}}
\newcommand*{\romanenumi}{\renewcommand*{\theenumi}{\roman{enumi}}\relabel}
\newcommand*{\Romanenumi}{\renewcommand*{\theenumi}{\Roman{enumi}}\relabel}
\newcommand*{\alphenumi}{\renewcommand*{\theenumi}{\alph{enumi}}\relabel}
\newcommand*{\Alphenumi}{\renewcommand*{\theenumi}{\Alph{enumi}}\relabel}

\let\oldtilde\tilde
\renewcommand*{\tilde}[1]{\mathchoice{\widetilde{#1}}{\widetilde{#1}}{\oldtilde{#1}}{\oldtilde{#1}}}
\let\oldhat\hat
\renewcommand*{\hat}[1]{\mathchoice{\widehat{#1}}{\widehat{#1}}{\oldhat{#1}}{\oldhat{#1}}}
\let\oldexists\exists
\renewcommand*{\exists}{\mathrel{\oldexists}}

\newcommand{\llbracket}{\mathopen{[\mkern -2.7mu[}}
\newcommand{\rrbracket}{\mathclose{]\mkern -2.7mu]}}
\newcommand{\llb}{\llbracket}
\newcommand{\rrb}{\rrbracket}

\providecommand{\zcref}[1]{\eqref{#1}}
\numberwithin{equation}{section}
\newtheorem{theo}{Theorem}[section]
\newtheorem{prop}[theo]{Proposition}
\newtheorem{coro}[theo]{Corollary}
\newtheorem{lemm}[theo]{Lemma}
\theoremstyle{definition}\newtheorem{defi}[theo]{Definition}
\newtheorem{rema}[theo]{Remark}
\title[One-sided minimizers for all means]{Simultaneous global inviscid Burgers flows with periodic Poisson forcing: original Theorem 1.4, proved as Proposition 2.5}
\author{Alexander Dunlap}
\date{Source: Annales Henri Lebesgue 8 (2025), 873--923; DOI 10.5802/ahl.250}
\begin{document}\maketitle
\noindent\textit{Editorial scope.} The counted unit is original Theorem 1.4, stated and proved as Proposition 2.5. The exact original setup, action functional and Definitions 1.1--1.3, the intro statement of Theorem 1.4, and the complete Section 2 core (concatenation, Proposition 2.1, Proposition 2.2 with its variational comparison, Definition 2.3 and Proposition 2.4) are retained as uncounted context. Later applications, the global-shock structure theorems, the Poisson verification and all figure inputs are omitted. Original published statement and display numbers are reproduced. Printed slips are kept literal without repair; every delivered body is byte-exact source text, with the publisher zcref macro served by an editorial wrapper alias that prints the same equation numbers as the published PDF.
\section{The forced inviscid Burgers equation and Lagrangian minimizers}
Let $\TT\coloneqq\RR/\ZZ$ and let $\pi\colon\RR\to\TT$ be the projection map. Let $\mu$ be a purely atomic measure on $ \RR\times\TT$ such that $\ssN\coloneqq\supp\mu$ is a discrete set. We will think of $\mu$ as a random measure, and our main example will be when $\mu$ is a realization of a homogeneous Poisson point process. We are interested in global solutions to the forced inviscid Burgers equation given by
\begin{subequations}\label{eq:uproblem}
\begin{align}
\partial_{t}u_{\theta}(t,x)+\partial_{x}\left(\frac{1}{2}u_{\theta}^{2}+\mu\right)(t,x) & =0,\qquad\theta,t\in\RR,x\in\TT;\label{eq:ueqn}
\\
\int_{\TT}u_{\theta}(t,x)\,\dif x & =\theta,\qquad\theta,t\in\RR.\label{eq:uintegral}
\end{align}
\end{subequations}

\begin{equation}\label{eq:Adef}
\clA_{\theta,s,t}[X]\coloneqq\frac{1}{2}\int_{s}^{t}(X'(r)-\theta)^{2}\,\dif r-\mu\bigl(\{(r,X(r))\st r\in[s,t)\}\bigr),
\end{equation}

\begin{defi}
Let $\Omega$ be the space of purely atomic measures $\mu$ on $\RR\times\TT$ such that $\ssN\coloneqq\supp\mu$ is a discrete set.
\end{defi}

\setcounter{equation}{7}
\begin{defi}[Lagrangian minimizers]\label{def:minimizers}
Let $\theta\in\RR$ and $s<t$.
\begin{enumerate}
\item We define the set $\scM_{s,y\mid t,x}^{\theta}$ comprising all paths $X\in H^{1}([s,t];\TT)$ with $X(s)=y$ and $X(t)=x$ such that if $Y$ is another such path, then $\clA_{\theta,s,t}[X]\le\clA_{\theta,s,t}[Y]$.
\item We define the set of \emph{one-sided minimizers}
\[
\scM_{t}^{\theta}\coloneqq\left\{X\in H_{\locc}^{1}((-\infty,t];\TT)\st X|_{[s,t]}\in\scM_{s,X(s)\mid t,X(t)}^{\theta}\text{ for all }s\le t\right\}.
\]
\item For $x\in\TT$, we define
\[
\scM_{t,x}^{\theta}\coloneqq\left\{X\in\scM_{t}^{\theta}\st X(t)=x\right\}.
\]
\item We define the partial order $\preccurlyeq$ on $\scM_{t,x}^{\theta}$ by
\begin{equation}\label{eq:partialorderdef}
X_{1}\preccurlyeq X_{2}\quad\text{if }\int_{r}^{t}X_{1}'(s)\,\dif s\ge\int_{r}^{t}X_{2}'(s)\,\dif s\quad\text{for all }r\le t.
\end{equation}
\end{enumerate}
\end{defi}

\begin{defi}\label{def:smallnoisezones}
We define $\widetilde{\Omega}_{1}$ to be the set of all $\mu\in\Omega$ such that there exists an $M=M(\mu)>0$ such that for all $t\in\RR$, there is an $s=s(t)\le t-M$ and a $y\in\TT$ such that
\begin{equation}\label{eq:intersectiononly}
\ssN\cap([s-M,s+M]\times\TT)=\{(s,y)\}
\end{equation}
and
\begin{equation}\label{eq:youshouldtake}
\mu(\{(s,y)\})>\frac{1}{4M},
\end{equation}
and moreover that we can choose $s(t)$ in such a way that
\begin{equation}\label{eq:stttoinfty}
\lim_{t\to+\infty}s(t)=+\infty.
\end{equation}
\end{defi}

\begin{theo}\label{thm:minimizers-exist}
Suppose that $\mu\in\widetilde{\Omega}_{1}$. For each $\theta,t\in\RR$ and $x\in\TT$, the set~$\scM_{t,x}^{\theta}$ is nonempty, consists of piecewise-linear paths connecting $(t,x)$ and points of $\ssN$, and has unique minimal and maximal elements $X_{\theta,t,x,\Left}$ and $X_{\theta,t,x,\Right}$, respectively, under the partial order $\preccurlyeq$.
\end{theo}

\section{Existence of one-sided minimizers}
In this section we prove Theorem~\ref{thm:minimizers-exist} on the existence of one-sided minimizers. First we must establish some basic properties of minimizers. If $X\in H^{1}([s,r];\TT)$ and $Y\in H^{1}([r,t];\TT)$ are such that $X(r)=Y(r)$, we define the concatenation $X\odot_{r}Y\colon[s,t]\to\TT$ by
\begin{equation}\label{eq:odotdef}
(X\odot_{r}Y)(q)\coloneqq
\begin{cases}
X(q), & \text{if }q\le r;\\
Y(q), & \text{if }q\ge r.
\end{cases}
\end{equation}

\begin{prop}\label{prop:basicproperties}
Suppose that $\mu\in\Omega$. Let $\theta\in\RR$, $-\infty<s<t<+\infty$, and $x,y\in\TT$.
\begin{enumerate}
\item \label{enu:straightlines}
The set $\scM_{s,y\mid t,x}^{\theta}$ is nonempty. Every $X\in\scM_{s,y\mid t,x}^{\theta}$ consists of straight line segments connecting points of $\{(t,x),(s,y)\}\cup\ssN$.
\item \label{enu:subpathalsominimizes}
If $\varnothing\ne[s',t']\subseteq[s,t]$ and $X\in\scM_{s,y\mid t,x}^{\theta}$, then $X|_{[s',t']}\in\scM_{s',X(s')\mid t',X(t')}^{\theta}$ as well.
\item \label{enu:pathsurgery}
If $r\in(s,t)$, $X\in\scM_{s,y\mid t,x}^{\theta}$ and $Y\in\scM_{s,y\mid r,X(r)}^{\theta}$, then $Y\odot_{r}X\in\scM_{s,y\mid t,x}^{\theta}$ as well.
\end{enumerate}
\end{prop}

\begin{proof}
The first point is a standard property of the convexity of the Dirichlet energy in~\eqref{eq:Adef}. For the second point, we note that if not, then we could modify $X$ on $[s',t']$ to improve the value of $\clA_{\theta,s,t}$, contradicting the assumption that $X\in\scM(\theta,s,t)$. To see the third point, note that
\[
\clA_{\theta,s,t}[Y\odot_{r}X]=\clA_{\theta,s,r}[Y]+\clA_{\theta,r,t}[X]\le\clA_{\theta,s,r}[X]+\clA_{\theta,s,r}[X]=\clA_{\theta,s,t}[X]
\]
by the definitions and part~\eqref{enu:subpathalsominimizes}.
\end{proof}

\begin{prop}\label{prop:gothroughy}
Suppose that $\mu\in\Omega$ and that $s\in\RR$, $M>0$, and $y\in\TT$ are such that~\eqref{eq:intersectiononly} and~\eqref{eq:youshouldtake} hold. Then, for any $\theta\in\RR$, $z_{1},z_{2}\in\TT$, and $X\in\scM_{s-M,z_{1}\mid s+M,z_{2}}^{\theta}$, we have $X(s)=y$.
\end{prop}

\begin{proof}
Assume towards a contradiction that we have some $\theta\in\RR$ and $X\in\scM_{s-M,z_{1}\mid s+M,z_{2}}^{\theta}$ such that $X(s)\ne y$. By the assumptions on $s$ and $y$, along with Proposition~\ref{prop:basicproperties}$\MK$\eqref{enu:straightlines}, we see that $X$ consists of a single straight line segment. Let
\begin{equation}\label{eq:xirange}
\xi\in(-1/2,1/2]
\end{equation}
be such that $X(s)+\xi=y$, and define
\[
Y(r)\coloneqq X(r)+\xi\cdot
\begin{cases}
\frac{s+M-r}{M} & \text{if }r\in[s,s+M];\\
\frac{r-s+M}{M} & \text{if }r\in[s-M,s];\\
0 & \text{otherwise}.
\end{cases}
\]
In particular, this means that $Y(s-M)=X(s-M)=z_{1}$, $Y(s+M)=X(s+M)=z_{2}$, and $Y(s)=X(s)+\xi=y$. Then we have
\begin{align*}
&\clA_{\theta,s-M,s+M}[Y]-\clA_{\theta,s-M,s+M}[X]\\
& \; =M\left[\frac{1}{2}\left(X'(s)-\theta+\xi/M\right)^{2}\!+\frac{1}{2}\left(X'(s)-\theta-\xi/M)^{2}-(X'(s)-\theta\right)^{2}\right]-\mu(\{(s,y)\})\\
& \; =\frac{\xi^{2}}{M}-\mu(\{(s,y)\})\overset{\zcref{eq:xirange}}{\le}\frac{1}{4M}-\mu(\{(s,y)\})\overset{\zcref{eq:youshouldtake}}{<}0.
\end{align*}
But since $Y(s\pm M)=X(s\pm M)$, this contradicts the assumption $X\in\scM_{s-M,z_{1}\mid s+M,z_{2}}^{\theta}$.
\end{proof}

\begin{defi}\label{def:Tstar}
Suppose that $\mu\in\widetilde{\Omega}_{1}$ and let $M(\mu)$ be as in Definition~\ref{def:smallnoisezones}. For $t\in\RR$, we define $T_{*}(t)$ to be the supremum of all $s<t$ such that there exists a $y\in\TT$ such that $(s,y)\in\ssN$ and $X(s)=y$ for all
\[
X\in\bigcup_{\substack{\theta\in\RR\\
x,z\in\TT } }\scM_{s-M(\mu),z\mid t,x}^{\theta}.
\]
\end{defi}

\setcounter{equation}{3}
\begin{prop}\label{prop:Tstartfinite}
Suppose that $\mu\in\widetilde{\Omega}_{1}$ and $t\in\RR$.
\begin{enumerate}
\item We have $T_{*}(t)>-\infty$.
\item There is a $y_{*}(t)\in\TT$ such that
\begin{equation}\label{eq:allpassthrough}
X(T_{*}(t))=y_{*}(t)\text{ for all }X\in\bigcup_{\substack{\theta\in\RR\\
x,z\in\TT } }\scM_{T_{*}(t)-M(\mu),z\mid t,x}^{\theta}.
\end{equation}
\item Finally, we have
\begin{equation}\label{eq:Tstarttoinfinity}
\lim_{t\to+\infty}T_{*}(t)=+\infty.
\end{equation}
\end{enumerate}
\end{prop}

\begin{proof}
Fix $t\in\RR$. Let $M>0$, $s\in(-\infty,t-M]$, and $y\in\TT$ be as in the definition of $\widetilde{\Omega}$. We claim that in fact $T_{*}(t)\ge s$. Proposition~\ref{prop:basicproperties}$\MK$\eqref{enu:subpathalsominimizes} implies that if $\theta\in\RR$, $x,z\in\TT$, and $X\in\scM_{s-M,z\mid t,x}^{\theta}$, then $X\in\scM_{s-M,z\mid s+M,X(s+M)}^{\theta}$ as well, and so by
Proposition~\ref{prop:Tstartfinite} we have $X(s)=y$. Since $(s,y)\in\ssN$ by definition, we have $T_{*}(t)\ge s>-\infty$, and the first assertion is proved. The second assertion is tantamount to asserting that the supremum in Definition~\ref{def:Tstar} is achieved; this follows from the discreteness of $\ssN$ and the compactness of $\TT$. Finally, \eqref{eq:Tstarttoinfinity} follows from~\eqref{eq:stttoinfty}.
\end{proof}

\begin{prop}\label{prop:globalminimizersexist}
Suppose that $\mu\in\widetilde{\Omega}_{1}$. For each $\theta,t\in\RR$ and $x\in\TT$, the set $\scM_{t,x}^{\theta}$ is nonempty. In particular, we have
\begin{equation}\label{eq:Mstarisrestriction}
\left\{X|_{[T_{*}(t),t]}\st X\in\scM_{t,x}^{\theta}\right\}=\scM_{T_{*}(t),y_{*}(t)\mid t,x}^{\theta}.
\end{equation}
Moreover, $\scM_{t,x}^{\theta}$ has minimal and maximal elements under the partial order $\preccurlyeq$.
\end{prop}

\begin{proof}
Since $\scM_{T_{*}(t),y_{*}(t)\mid t,x}^{\theta}$ is nonempty by
Proposition~\ref{prop:basicproperties}$\MK$\eqref{enu:straightlines}, to prove that the set $\scM_{t,x}^{\theta}$ is nonempty it suffices to prove~\eqref{eq:Mstarisrestriction}. The ``$\subseteq$'' direction is an immediate consequence of Proposition~\ref{prop:basicproperties}$\MK$\eqref{enu:subpathalsominimizes} and the definition of $\scM_{t,x}^{\theta}$, so we turn our attention to the ``$\supseteq$'' direction. In other words, given $X_{1}\in\scM_{T_{*}(t),y_{*}(t)\mid t,x}^{\theta}$, we seek to extend $X_{1}$ to an element $X$ of $\scM_{t,x}^{\theta}$. Let $t_{0}=t$, and let $t_{k}=T_{*}(t_{k-1})$ and $y_{k}=y_{*}(t_{k-1})$ for $k\ge1$. We note that, since $t_{k}\le t_{k-1}-M(\mu)$, we have
\begin{equation}\label{eq:tksfillup}
\bigsqcup_{k=1}^{\infty}(t_{k},t_{k-1}]=(-\infty,t].
\end{equation}
For $k\ge2$, let $X_{k}\in\scM_{t_{k},y_{k}\mid t_{k-1},y_{k-1}}^{\theta}$. (This inclusion is satisfied for $k=1$ as well, but $X_{1}$ has already been chosen.) Now for $s\in[t_{k},t_{k-1})$, define $X(s)\coloneqq X_{k}(s)$, so $X$ is defined as an element of $H_{\locc}^{1}((-\infty,t];\TT)$ by~\eqref{eq:tksfillup} and the fact that $X_{k}(t_{k-1})=y_{k-1}=X_{k-1}(t_{k-1})$ by definition.

We claim that $X\in\scM_{t,x}^{\theta}$. Let $s<t$ and let $k$ be large enough that $t_{k}\le s$. Suppose that $z\in\TT$ and that $Y\in\scM_{t_{k}-M,z\mid t,x}^{\theta}$. Then, by~\eqref{eq:allpassthrough}, we have $Y(t_{j})=y_{j}=X(t_{j})$ whenever $j\le k$. This means that $Y|_{[t_{j},t_{j-1}]}\in\scM_{t_{k},y_{k}\mid t_{k-1},y_{k-1}}^{\theta}$. Since the same is true for $X|_{[t_{j},t_{j-1}]}$, we have $\clA_{\theta,t_{j},t_{j-1}}[X]=\clA_{\theta,t_{j},t_{j-1}}[Y]$. Summing this up over all $j$, we obtain
\begin{equation}\label{eq:YXsame}
\clA_{\theta,t_{k},t}[Y]=\clA_{\theta,t_{k},t}[X].
\end{equation}
Now since $Y\in\scM_{t_{k},y_{k}\mid t,x}^{\theta}$ by Proposition~\ref{prop:basicproperties}$\MK$\eqref{enu:subpathalsominimizes},
\eqref{eq:YXsame} means that $X\in\scM_{t_{k},y_{k}\mid t,x}^{\theta}$ as well. But since $t_{k}\le s$, we can apply Proposition~\ref{prop:basicproperties}$\MK$\eqref{enu:subpathalsominimizes} once again to see that $X\in\scM_{s,X(s)\mid t,x}^{\theta}$ Since this is true for any $s<t$, we conclude that $X\in\scM_{t,x}^{\theta}$.

To show that $\scM_{t,x}^{\theta}$ has minimal/leftmost and maximal/rightmost elements under~$\preccurlyeq$, we observe that each $\scM_{t_{k},y_{k}\mid t_{k-1},y_{k-1}}^{\theta}$ is finite and has leftmost and rightmost elements (as can be seen using Proposition~\ref{prop:basicproperties}$\MK$\eqref{enu:pathsurgery} to build a path that is weakly to the left/right of any other minimizer). Then we note that a concatenation of these leftmost and rightmost elements in a similar manner to the above argument will yield leftmost and rightmost elements of $\scM_{t,x}^{\theta}$.
\end{proof}

\end{document}
